\documentclass[11pt, a4paper]{article}

\usepackage[margin=2.5cm]{geometry}
\usepackage{amsmath, amssymb, amsthm}
\usepackage{booktabs}
\usepackage{hyperref}
\setlength{\parskip}{0.5em}
\setlength{\parindent}{0pt}

\hypersetup{
  colorlinks = true,
  linkcolor  = blue,
  urlcolor   = blue
}

\newtheorem{remark}{Remark}
\newtheorem{principle}{Design Principle}

\title{\textbf{Application Module: Design Document}\\[0.4em]
  \large Large-Scale Panel Data Analysis via Online TWFE\\[0.2em]
  \normalsize Hwang \& Lee (2026) ---
  \textit{Online Updating for Linear Panel Regressions}}
\author{}
\date{March 2026}

\begin{document}
\maketitle
\tableofcontents
\bigskip\bigskip

% =============================================================================
\section{Overview and Motivation}
% =============================================================================

The simulation study (Section~5.2) demonstrates that the online TWFE algorithm
achieves numerical equivalence with a full-sample batch estimator while
processing units sequentially.  In that setting, data are \emph{generated}
unit by unit, so the full panel is never materialised in memory.

The application module addresses a complementary and practically important
scenario: a \emph{pre-existing} large dataset---stored as a flat file or in a
database---that is too large to load into RAM at once.  The canonical example
is a dataset with $N = 10^8$ individuals and $T = 10$ time periods, whose raw
size may easily exceed 10~GB.

Standard batch estimators such as \texttt{plm} require the entire
\texttt{pdata.frame} to reside in memory before estimation begins.  The
benchmark results already show that \texttt{plm} exhausts a 16~GB memory limit
at $N \approx 2 \times 10^7$.  The online TWFE algorithm, by contrast,
maintains a fixed-size state object of $O(p^2)$ scalars regardless of $N$,
and processes the data one chunk at a time.

\begin{remark}
The distinction between the simulation and application settings lies entirely
in \emph{how data enter the algorithm}, not in the algorithm itself.  In both
cases the same algebraically exact update equations (Propositions~1--3) are
applied; only the data ingestion pipeline differs.
\end{remark}

% =============================================================================
\section{Model and Estimands}
% =============================================================================

The target model is the two-way fixed-effects (TWFE) panel regression
\[
  Y_{it} \;=\; X_{it}'\beta + \alpha_i + \lambda_t + u_{it},
  \qquad i = 1,\ldots,N,\quad t \in \mathcal{S}_i \subseteq \{1,\ldots,T\},
\]
where $\alpha_i$ denotes individual fixed effects and $\lambda_t$ calendar-time
fixed effects.  The panel is \emph{unbalanced}: each unit $i$ is observed on a
subset $\mathcal{S}_i$ of the $T$ calendar periods.

The three estimands of interest are:
\begin{enumerate}
  \item[(i)] the slope coefficient vector $\hat\beta$ (via Algorithm~1);
  \item[(ii)] the classical variance $\widehat{V}_0 = \hat\sigma^2
        (\dot{Z}'\dot{Z})^{-1}$; and
  \item[(iii)] the cluster-robust (Arellano HC0) variance $\hat{V}_{cr}$.
\end{enumerate}
All three are produced in a \emph{single pass} over the data chunks.

% =============================================================================
\section{Prerequisites and Design Constraints}
% =============================================================================

\subsection{Unit-Boundary Constraint}

The within transformation for individual $i$,
\[
  \dot{Z}_i = Z_i - \bar{Z}_i, \qquad \dot{Y}_i = Y_i - \bar{Y}_i,
\]
requires \emph{all} observations of unit $i$ to be available simultaneously.
Consequently, a fundamental constraint on any chunked design is:

\begin{principle}[Unit atomicity]
No chunk may split the observations of a single unit across two chunks.
Every unit's rows must appear entirely within one chunk.
\end{principle}

This implies that chunks are defined in terms of \emph{number of units}, not
number of rows.  For an unbalanced panel, the row count of a chunk of $C$ units
is $\sum_{i \in \mathrm{chunk}} T_i$, which varies across chunks.

\subsection{Sorted Input Requirement}

Unit atomicity can only be guaranteed if the dataset is sorted by the
individual identifier (\texttt{id}) prior to chunked reading.

\begin{principle}[Sorted input]
The dataset must be sorted by \texttt{id} in non-decreasing order before the
chunked analysis begins.
\end{principle}

A dataset sorted by calendar time, for instance, would interleave the rows of
different units and make unit-boundary detection impossible without a full pass.

\subsection{Pre-specified Calendar Time Support}

Let $\mathcal{T} = \{t_1, t_2, \ldots, t_T\}$ denote the full set of calendar
times present in the dataset.  The $p = k + T - 1$ dimensional parameter
space is fixed at initialisation.  If $\mathcal{T}$ were allowed to grow
during chunked processing, Algorithm~3 (new calendar time) would be triggered,
expanding the parameter dimension and requiring re-projection of previously
accumulated sufficient statistics onto the enlarged parameter space.

\begin{principle}[Pre-specified $T_{\mathrm{support}}$]
The full calendar time support $T$ is determined before processing begins and
passed to the initialisation step.  All chunks are assumed to contain only
time values within $\{1,\ldots,T\}$.
\end{principle}

In practice, $T$ is either known from the study design (e.g.\ a 10-year panel
has $T = 10$) or is determined with a single lightweight scan of only the
\texttt{time} column of the file, which is fast even for very large files.

% =============================================================================
\section{Four-Phase Workflow}
% =============================================================================

\subsection{Phase 0 --- Data Inspection}

Before any estimation, the analyst performs a lightweight scan to collect
structural metadata:
\begin{itemize}
  \item Column names and data types (\texttt{id}, \texttt{time}, \texttt{y},
        $x_1,\ldots,x_k$).
  \item The set of unique calendar times $\{t_1,\ldots,t_T\}$, from which
        $T_{\mathrm{support}}$ and \texttt{baseline\_time} are determined.
  \item Whether the file is already sorted by \texttt{id}.
\end{itemize}

This phase reads only the \texttt{id} and \texttt{time} columns (a small
fraction of the total file size) and produces no large intermediate objects.

\subsection{Phase 1 --- Preprocessing: Sort}

If the file is not already sorted by \texttt{id}, an external sort is required.
For a 10~GB file this is infeasible with R's in-memory sort; instead, one of
the following approaches is used:
\begin{itemize}
  \item Unix \texttt{sort} utility (fastest; operates directly on disk).
  \item Apache Arrow (\texttt{open\_dataset} $\to$ \texttt{arrange} $\to$
        \texttt{write\_csv\_arrow}).
  \item A database \texttt{ORDER BY id} query.
\end{itemize}
Phase~1 is a \emph{one-time cost}.  If the data source is a database or a
file system that supports querying by id range, this phase may be bypassed.

\subsection{Phase 2 --- Chunked State Accumulation}

This is the core of the application module.  Phase~2 implements the pipeline:
\[
  \texttt{otwfe\_init} \;\to\;
  \underbrace{\texttt{otwfe\_update}(\text{chunk 1})
    \;\to\; \cdots \;\to\;
    \texttt{otwfe\_update}(\text{chunk } C)}_{\text{repeat for each chunk}}.
\]

\paragraph{Chunk boundary detection.}
Given the sorted file, the unit boundaries are identified by run-length
encoding of the \texttt{id} column.  This produces, for each unit, the number
of rows it occupies.  Cumulative sums of these lengths yield exact row offsets
for every unit, allowing each chunk to be read as a contiguous block of exactly
$C$ units.  This scan is performed once and the boundary table (a vector of
$N$ integers) is retained in memory; for $N = 10^8$, this amounts to
approximately 400~MB---a fixed cost independent of the file size.

\paragraph{Initialisation (\texttt{otwfe\_init}).}
The handle is created with:
\begin{itemize}
  \item the covariate column names $x_1,\ldots,x_k$;
  \item $T_{\mathrm{support}}$ (from Phase~0);
  \item \texttt{baseline\_time} (the omitted time dummy); and
  \item a verbosity flag.
\end{itemize}
At this point no data have been read.  The state object $\mathcal{S}$ is
uninitialised; the handle is an empty shell.

\paragraph{First chunk (\texttt{otwfe\_update}, call 1).}
Processing the first chunk involves two sub-steps:

\begin{enumerate}
  \item \textbf{Warm-up initialisation.}  A small subset of units (at least
  $3p$ units, minimum 30) is selected from the first chunk by a greedy cover
  algorithm that guarantees every calendar time $t \in \{1,\ldots,T\}$ is
  represented at least once in the warm-up set.  These units are used to
  compute an initial OLS estimate via the full within-regression, producing the
  initial state
  \[
    \mathcal{S}_0 = \bigl\{
      (\dot{Z}'\dot{Z})^{-1},\;
      \hat\beta,\;
      \hat\sigma^2,\;
      N_{\mathrm{warm}},\;
      n_{\mathrm{warm}},\;
      M_{ss},\;
      A_N,\;
      B_N
    \bigr\}.
  \]

  \item \textbf{Remaining units.}  All units in the first chunk beyond the
  warm-up set are processed by the Rcpp batch update (\texttt{alg1\_batch\_cpp}),
  which adds their sufficient statistics to $\mathcal{S}$ in a single pass.
\end{enumerate}

\paragraph{Subsequent chunks (\texttt{otwfe\_update}, calls 2 onward).}
Each subsequent chunk contains only new units not yet seen by the algorithm.
The Rcpp batch routine processes the entire chunk in a single pass:
\begin{enumerate}
  \item For each unit $i$ in the chunk, compute
        $S_i = \dot{Z}_i'\dot{Z}_i$,\; $s_i = \dot{Z}_i'\dot{Y}_i$,\;
        $\mathrm{SSy}_i = \dot{Y}_i'\dot{Y}_i$, and $\mathrm{vec}(S_i)$.
  \item Accumulate $\sum S_i$, $\sum s_i$, $\sum \mathrm{SSy}_i$
        into running totals.
  \item Accumulate the three Vcr building blocks:
  \[
    M_{ss} \mathrel{+}= s_i s_i', \qquad
    A_N    \mathrel{+}= s_i \operatorname{vec}(S_i)', \qquad
    B_N    \mathrel{+}= \operatorname{vec}(S_i)\operatorname{vec}(S_i)'.
  \]
  \item Update $(\dot{Z}'\dot{Z})^{-1}$, $\hat\beta$, and $\hat\sigma^2$
        using the rank-update formula.
\end{enumerate}

The key property that makes this chunk-safe is:

\begin{remark}[Chunk-safe Vcr accumulation]
The matrices $M_{ss}$, $A_N$, $B_N$ are defined as sums over individual units.
Because addition is commutative and associative, these sums can be split across
any number of chunks and accumulated in any order.  The final Vcr is computed
\emph{after all chunks have been processed}, using the final $\hat\beta$:
\[
  M \;=\;
    M_{ss}
    - A_N\bigl(\hat\beta' \otimes I_p\bigr)
    - \bigl(\hat\beta' \otimes I_p\bigr)' A_N'
    + \bigl(\hat\beta' \otimes I_p\bigr) B_N \bigl(\hat\beta' \otimes I_p\bigr)',
\]
\[
  \hat{V}_{cr} \;=\; (\dot{Z}'\dot{Z})^{-1}\, M\, (\dot{Z}'\dot{Z})^{-1}.
\]
This formula is algebraically exact regardless of how many chunks were used or
in what order units were processed.
\end{remark}

\paragraph{Memory profile.}
After each call to \texttt{otwfe\_update}, the chunk data frame is immediately
deleted from memory.  The only object persisting between chunks is the state
handle $\mathcal{S}$, whose size is:
\[
  \underbrace{p \times p}_{\text{inv}} +
  \underbrace{p}_{\hat\beta} +
  \underbrace{1}_{\hat\sigma^2} +
  \underbrace{p \times p}_{M_{ss}} +
  \underbrace{p \times p^2}_{A_N} +
  \underbrace{p^2 \times p^2}_{B_N}
  \;=\; O(p^4) \text{ scalars.}
\]
For $p = 4$ ($k=2$, $T=3$), this is fewer than 200 double-precision numbers
--- well under one kilobyte.  This footprint is \emph{independent of $N$}.

\subsection{Phase 3 --- Finalisation}

After all chunks have been processed, \texttt{otwfe\_finalize} performs:
\begin{enumerate}
  \item Compute $\hat{V}_{cr}$ from the accumulated $M_{ss}$, $A_N$, $B_N$,
        and the final $\hat\beta$ using the formula in Remark~2.
  \item Package $\hat\beta$, $\hat{V}_0$, $\hat{V}_{cr}$, and metadata into
        the standard \texttt{otwfe} return object.
\end{enumerate}
This step is $O(p^4)$ and takes negligible time.

% =============================================================================
\section{Correctness Guarantees}
% =============================================================================

The chunked pipeline is \emph{not} an approximation.  The following
equivalences hold exactly (up to floating-point rounding at double precision):

\begin{enumerate}
  \item[(i)] $\hat\beta_{\mathrm{chunk}} = \hat\beta_{\mathrm{batch}}$:
        The coefficient vector produced by processing $C$-unit chunks is
        identical to the OLS estimate from a single batch regression on the
        full dataset.

  \item[(ii)] $\hat{V}_{0,\mathrm{chunk}} = \hat{V}_{0,\mathrm{batch}}$:
        The classical variance $\hat\sigma^2 (\dot{Z}'\dot{Z})^{-1}$ is exact
        because both $\hat\sigma^2$ and $(\dot{Z}'\dot{Z})^{-1}$ accumulate
        exactly across chunks.

  \item[(iii)] $\hat{V}_{cr,\mathrm{chunk}} = \hat{V}_{cr,\mathrm{batch}}$:
        By Remark~2, the single-pass Vcr formula evaluated at the final
        $\hat\beta$ yields the same result as the Arellano HC0 sandwich
        estimator applied to the full sample.
\end{enumerate}

These equivalences were verified numerically in the simulation study at
$N = 10^6$, $5 \times 10^6$, and $10^7$, with maximum coefficient discrepancies
on the order of $10^{-13}$, consistent with double-precision arithmetic.

% =============================================================================
\section{Comparison with Batch Methods}
% =============================================================================

\begin{table}[h]
  \centering
  \caption{Online TWFE (chunk API) vs.\ \texttt{plm} (batch)}
  \begin{tabular}{lcc}
    \toprule
    Property & Online TWFE (chunk) & \texttt{plm} \\
    \midrule
    Full dataset in RAM required      & No  & Yes \\
    RAM usage scales with $N$         & No ($O(p^4)$) & Yes ($O(nk)$) \\
    Feasible at $N = 2 \times 10^7$   & Yes & No (OOM on 16\,GB) \\
    Algebraically exact               & Yes & Yes \\
    Cluster-robust SE (Arellano HC0)  & Yes (single pass) & Yes (post-estimation) \\
    Requires sorted input             & Yes & No \\
    One-time preprocessing cost       & Yes (sort) & None \\
    \bottomrule
  \end{tabular}
\end{table}

The primary trade-off is that the chunk API requires the dataset to be sorted
by \texttt{id} in advance.  For a one-time analysis of a persistent dataset,
this is a mild requirement.  For a truly streaming environment where data are
generated in unit order (as in the simulation study), sorting is not needed.

% =============================================================================
\section{Implementation Plan}
% =============================================================================

\subsection{Files to be Created or Extended}

\textbf{\texttt{Simulation/sim\_dgp.R} (extension):}
A new function \texttt{make\_panel\_52\_stream()} will be added to the
existing DGP file.  This function returns a \emph{closure-based iterator} that,
on each call to \texttt{next\_chunk()}, generates and returns exactly one chunk
of $C$ units without retaining any state beyond the chunk index, an id offset
counter, and the pre-generated $\lambda_t$ vector ($T$ scalars).

\textbf{\texttt{Application/streaming\_bigdata.R}:}
Main demonstration script that exercises the four-phase workflow on datasets of
increasing size, reports timing and accuracy, and produces the summary table
for the paper.

\subsection{Streaming DGP Design (\texttt{make\_panel\_52\_stream})}

The function is designed with the following properties:

\begin{itemize}
  \item \textbf{Minimal pre-allocation.}  At construction, only $\lambda_t$
        ($T$ scalars) is drawn.  Nothing proportional to $N$ is allocated.

  \item \textbf{Per-chunk generation.}  Each call to \texttt{next\_chunk()}
        draws $\alpha_i$, $T_i$, calendar time subsets, $X_{it}$, $u_{it}$,
        and $Y_{it}$ for the next $C$ units, constructs the data frame, and
        returns it.  The generator's internal state advances by $C$ units.

  \item \textbf{Reproducibility.}  Each chunk uses a seed derived from the
        global seed and the chunk index.  This ensures that any particular
        chunk produces the same data regardless of how many chunks have been
        processed before it.

  \item \textbf{Globally unique IDs.}  Chunk $c$ produces unit IDs from
        $(c-1)C + 1$ to $\min(cC,\, N_{\mathrm{total}})$, so IDs never
        repeat and the data are naturally sorted by \texttt{id} within each
        chunk and across chunks.

  \item \textbf{Termination.}  When all $N_{\mathrm{total}}$ units have been
        generated, \texttt{next\_chunk()} returns \texttt{NULL}, signalling
        the end of the stream.
\end{itemize}

\subsection{Application Script Structure (\texttt{streaming\_bigdata.R})}

The script is organised into three sections:

\textbf{Section 1 --- Correctness verification (small $N$).}
For $N = 1{,}000$ and $T = 3$, run both the streaming pipeline and
\texttt{plm}, and confirm that the maximum absolute differences
\[
  \|\hat\beta_{\mathrm{online}} - \hat\beta_{\mathrm{plm}}\|_\infty,\quad
  \|\hat{V}_0^{\mathrm{online}} - \hat{V}_0^{\mathrm{plm}}\|_\infty,\quad
  \|\hat{V}_{cr}^{\mathrm{online}} - \hat{V}_{cr}^{\mathrm{plm}}\|_\infty
\]
are all at machine-precision level ($\lesssim 10^{-13}$).

\textbf{Section 2 --- Scalability benchmark.}
For $N \in \{10^6,\, 10^7,\, 10^8,\, 5 \times 10^8\}$ (or as large as time
permits), run the streaming pipeline and record:
\begin{itemize}
  \item total wall-clock time (DGP generation + update + finalize);
  \item peak per-chunk memory (size of one chunk data frame in MB);
  \item state handle size (constant across $N$, reported once);
  \item estimation bias $\|\hat\beta - \beta_{\mathrm{true}}\|_\infty$.
\end{itemize}

\textbf{Section 3 --- Summary table.}
A formatted table comparing all runs, with a note that \texttt{plm} fails
beyond $N \approx 2 \times 10^7$ on a 16~GB machine (from the simulation
benchmark).

% =============================================================================
\section{Known Limitations and Scope}
% =============================================================================

\begin{itemize}
  \item \textbf{Algorithm~3 not used.}  The application module assumes a
        pre-specified and fixed $T_{\mathrm{support}}$.  If new calendar times
        were to appear in later chunks, the parameter dimension would need to
        expand (Algorithm~3), which is not implemented in the current chunk API
        and is left for future work.

  \item \textbf{New observations for existing units not handled.}
        The current pipeline assumes all units are new (Algorithm~1 only).
        Algorithm~2 (new time observation for an existing unit) is not
        implemented in the chunk API.

  \item \textbf{Sorting cost.}  For very large files, the external sort in
        Phase~1 may require significant time and temporary disk space.
        Datasets originating from a database can bypass this by issuing an
        \texttt{ORDER BY id} query at read time.

  \item \textbf{Boundary scan memory.}  The run-length encoding of the
        \texttt{id} column produces a vector of $N$ integers ($\sim$400~MB for
        $N = 10^8$).  For $N \geq 10^9$, this may itself become a memory
        concern.  A future improvement would read chunks adaptively without a
        pre-computed boundary table.
\end{itemize}

\end{document}
